

%% ===================================================================
\section{RTO의 정적 한계}
\label{sec:lit-rto-static}
%% ===================================================================

\chg{앞 절에서 살펴본 BCM의 핵심 산출물 가운데 하나가 RTO이다. 그러나 현행 RTO는 복구 환경의 변화와 무관하게 단일 정적 값으로 고정된다는 구조적 제약을 안고 있다. 이 절에서는 그 한계를 정적 설정의 구조적 문제와 선행연구의 접근을 중심으로 검토한다.}

%% -------------------------------------------------------------------
\subsection{정적 RTO 설정의 구조적 문제}
\label{subsec:static-rto-problem}
%% -------------------------------------------------------------------

전통적 RTO 설정 방식에는 세 가지 근본적 한계가 존재하며, 이는 본 논문
\secref{sec:problem_statement}(연구 필요성)에서 제시한 세 가지 연구 필요성과 직접 대응된다.

\chg{첫째는 관측성 미반영이다. 정적 RTO는 BIA에서 도출된 고정값으로서 조직의 운영
가시성, 곧 관측성(observability)을 반영하지 못한다. 시스템 복잡도와 인력 역량, 기술 환경 등
운영 환경이 시간과 상황에 따라 변화하는 현실에서 관측성이 높은 시스템은 장애 원인을 신속히
파악하여 복구 시간을 단축할 수 있으나, 전통적 BIA는 이러한 운영 역량 차이를 RTO에
반영하는 메커니즘을 갖추지 못한다. 그 결과 관측성이 높은 조건에서는 복구 자원의 과다
예비(over-provisioning)가, 낮은 조건에서는 과소 예비(under-provisioning)가 동시에 초래된다.
둘째는 일상적 불확실성의 정량화 부재이다. 복구 과정에서 발생하는 장비 가용성과 인력 동원 속도,
외부 공급망 의존도 같은 일상적 불확실성이 정량적으로 모델링되지 않는다. 실제 복구 시간은 다양한
예측 불가능한 요인에 의해 목표 시간을 초과할 수 있으나, 기존 접근법은 이러한 정규분포(Gaussian)적
변동을 체계적으로 반영하지 못한다. 셋째는 극단 사건의 미모형화이다. 사이버공격에 의한 데이터
센터 전면 장애나 복합 재난으로 인한 공급망 동시 중단처럼 일상적 변동을 넘는 극단적 사건에 의한
복구 지연은, 정규분포 가정으로는 포착할 수 없는 파레토(Pareto, 중꼬리, heavy-tail) 분포 특성을
보인다~\citep{coles2001extremes, embrechts1997extremal, taleb2007blackswan}.}


%% -------------------------------------------------------------------
\subsection{선행 연구에서의 RTO 접근}
\label{subsec:prior-rto-approaches}
%% -------------------------------------------------------------------

기존의 RTO 관련 연구는 대체로 BIA 방법론의 개선, RTO 달성을 위한 기술적 수단
(이중화, 백업 전략), 또는 특정 산업 분야에서의 사례 분석에 집중하고 있다.
\citet{snedaker2013bcdr}는 IT 재해복구와 업무연속성 계획 수립의 체계적 방법론을
제시하였으나, RTO 자체를 동적으로 조정하는 수리 모델은 다루지 않았다.
\citet{cerullo2004bcp}는 금융 서비스 산업에서의 BCP 수립 사례를 분석하여
RTO의 실무적 중요성을 강조하였으나, RTO의 동적 특성에 대한 논의는 부재하였다.

\chg{최근에는 업무연속성과 재해복구를 통합하여 조직 복원력 관점에서 정량적으로 접근하려는 시도가 이어지고 있다. \citet{sahebjamnia2015integrated}는 핵심 업무의 재개와 복구를 동시에 고려하는 다목적 혼합정수계획 모형으로 복구 자원의 최적 배분 문제를 정식화하였고, \citet{torabi2016enhanced}는 최선-최악 방법(Best--Worst Method)에 기반한 향상된 위험평가 체계를 제안하여 BCM의 위험 정량화 수준을 끌어올렸다. 그러나 이들 정량적 프레임워크에서도 RTO 자체는 위험평가와 자원배분의 고정된 입력 매개변수로 취급되며, 운영 환경의 관측성이나 복구 시간의 확률적 변동에 따라 RTO를 동적으로 조정하는 메커니즘은 여전히 다루어지지 않는다.}

\citet{cheung2019quickhit}, \citet{jang2021bcms} 그리고 \citet{jung2023bcms}을 포함한 국내 BCM
연구 또한 표준 적용과 조직 성과 평가에 집중하고 있어, RTO의 동적 조정이나
운영 조건에 따른 적응적 복구 모델에 대한 연구가 상대적으로 부재함을 드러낸다.
\citet{lee2026drto}이 관측성 채널(C0$\to$C1)을 도입하여
초기 접근을 수행하였으나, 불확실성 채널(C2/C3)의 체계적 통합은 향후 과제로
남아있었다.
이러한 세 가지 한계와 향후과제가 본 연구에서 $\DRTO$ 프레임워크를 통해 해소하고자 하는
핵심 문제의식의 출발점이다.

\chg{이러한 적응적 복구를 실현하려면 복구 환경의 상태를 실시간으로 관측할 수 있어야 하며, 다음 절에서는 그 이론적 토대인 관측성 개념을 검토한다.}
